\documentclass[10pt]{article}

\usepackage[utf8]{inputenc}
\usepackage[english]{babel}
\usepackage[margin=0.85in]{geometry}
\usepackage{amsmath,amsthm,amssymb,stmaryrd}
\usepackage{hyperref,thmtools,enumitem,xcolor}

\hypersetup{
    colorlinks,
    linkcolor={red!50!black},
    citecolor={blue!50!black},
    urlcolor={blue!80!black}
}

\declaretheorem[name=Theorem,numberwithin=section]{theorem}
\declaretheorem[name=Lemma,sibling=theorem]{lemma}
\declaretheorem[name=Proposition,sibling=theorem]{proposition}
\declaretheorem[name=Corollary,sibling=theorem]{corollary}
\declaretheorem[name=Definition,sibling=theorem,style=definition]{definition}
\declaretheorem[name=Remark,sibling=theorem,style=remark]{remark}
\setlist[itemize]{parsep=0pt,topsep=2pt,before=\leavevmode}
\setlist[enumerate]{parsep=0pt,topsep=2pt,before=\leavevmode}
\setlength{\parindent}{1em}
\setlength{\parskip}{.15em}

\newcommand{\N}{\mathbb{N}}
\newcommand{\Asm}{\mathbf{Asm}}
\newcommand{\Prop}{\mathsf{Prop}}
\newcommand{\Type}{\mathsf{Type}}
\newcommand{\Set}{\mathsf{Set}}
\newcommand{\CIC}{\mathrm{CIC}}
\newcommand{\ZF}{\mathrm{ZF}}
\newcommand{\ZFC}{\mathrm{ZFC}}
\newcommand{\IZF}{\mathrm{IZF}}
\newcommand{\HAH}{\mathrm{HAH}}
\newcommand{\EM}{\mathrm{EM}}
\newcommand{\Con}{\mathrm{Con}}
\newcommand{\code}[1]{\ulcorner #1\urcorner}
\newcommand{\El}{\mathrm{El}}
\newcommand{\dom}{\operatorname{dom}}
\newcommand{\rk}{\operatorname{rank}}
\newcommand{\rz}{\Vdash}
\newcommand{\ap}{\cdot}
\newcommand{\U}{\mathcal{U}}
\newcommand{\PP}{\mathcal{P}}
\newcommand{\Nat}{\mathsf{nat}}
\newcommand{\Acc}{\mathsf{Acc}}
\newcommand{\uchoice}{\mathsf{uchoice}}
\newcommand{\CT}{\mathrm{CT}}
\newcommand{\sem}[1]{\llbracket #1\rrbracket}

\begin{document}

\title{Universes Without Inaccessibles:\\
  \large a ZF model of the Calculus of Inductive Constructions with $\omega$ universes}
\author{Mario Carneiro}
% Drafted with Claude; adjust authorship/acknowledgement as you see fit.
\date{Draft, \today}
\maketitle

\noindent\fbox{\parbox{\dimexpr\linewidth-2\fboxsep}{\textbf{Status (2026-09-18).}
Lemma~\ref{lem:lfp}(2) is false and Theorem~\ref{thm:main} is withdrawn:
Proposition~\ref{prop:noset} shows that \emph{no} set-sized universe of
the kind constructed in Section~\ref{sec:model} exists unless an
inaccessible cardinal does. Sections~\ref{sec:theories}--\ref{sec:lower},
Remark~\ref{rem:diff}(iii), Proposition~\ref{prop:validates} for the
one-universe model, and Remark~\ref{rem:prop} stand. The abstract and
introduction describe the intended result and are left as written for now.}}

\begin{abstract}
The standard set-theoretic models of the Calculus of Inductive Constructions
interpret each universe $\Type_i$ as a stage $V_{\kappa_i}$ of the cumulative
hierarchy and so assume one inaccessible cardinal per universe. We show these
cardinals are an artifact of the model. Building on the assembly model of van
den Berg and Otten, we interpret each universe as the \emph{least} set of
codes closed under the type formers, where a family of types is a computably
tracked map into the universe. A tracker determines the computational
skeleton of a family, so the closure is reached after countably many steps
and the whole hierarchy is a set of rank $\omega_1$, although the universes
themselves are large. We obtain
$\ZF\vdash\Con(\CIC_\omega+\EM)$ with $\EM$ excluded middle in $\Prop$, using
no choice and no large cardinals. The model refutes unique choice, exactly
the principle Werner and Barras needed to interpret Replacement, so unique
choice rather than the universes is what separates $\CIC$ from $\ZF$.
Conversely, $\CIC$ with two universes proves the consistency of higher-order
Heyting arithmetic and of its transfinite extensions.
\end{abstract}

\section{Introduction}

How strong is the Calculus of Inductive Constructions with a cumulative
hierarchy $\Prop\subseteq\Type_0\subseteq\Type_1\subseteq\cdots$? The usual
answer comes from the sets-in-types-and-types-in-sets models of Werner
\cite{werner97}, Barras \cite{barras10}, and Timany and Sozeau
\cite{timanysozeau17}: $\Type_i$ is interpreted as $V_{\kappa_i}$ with
$\kappa_i$ inaccessible, so $\ZFC$ plus $n$ inaccessibles proves
$\Con(\CIC_n)$. The inaccessibles are forced by the model, not by the theory.
$V_\kappa$ is closed under products of \emph{all} set-theoretic families
exactly when $\kappa$ is inaccessible, and in that model a family is any set
function; but the type theory can only build families it can name, and it
cannot even prove Replacement for its own sets-as-trees: Werner needed a
``type-theoretic description axiom'' and Barras a unique choice axiom
$\uchoice$, and even then Barras ``could not prove the collection axiom''
nor closure of Grothendieck universes under functional replacement
\cite{barras10}. Timany and Sozeau state the matching lower bound only as a
conjecture.

This note gives a model with no inaccessible: the universes are taken as
small as the type theory can force them to be. We work in the category
$\Asm$ of assemblies over Kleene's first algebra, as van den Berg and Otten
\cite{vdbergotten23} do for one universe, and interpret $\Type_i$ as a
\emph{Tarski universe}: a set of codes with a decoding into $\Asm$,
obtained as the least fixed point of a monotone closure operator. A family
of types over $A$ is a morphism from $A$ into the universe, tracked by a
natural number that determines the family up to its propositional parts.
The stage at which a code enters the universe depends only on this
computational skeleton, so the closure stabilises after countably many
stages even though the universes contain continuum many types. Nothing has
to be regular; only the existence of each product is used.

\paragraph{Main results.}
Let $\CIC_n$ be the calculus with sorts $\Prop$ and $\Type_0,\dots,\Type_{n-1}$
and $\CIC_\omega$ the union (Section~\ref{sec:theories}); let $\EM$ be
excluded middle in $\Prop$.
\begin{itemize}
\item $\ZF\vdash\Con(\CIC_\omega+\EM)$ (Theorem~\ref{thm:main}). The model is
  a set of rank $\omega_1$, and Replacement is used only for a transfinite
  recursion of countable length and for the sets $\beth_\alpha$,
  $\alpha<\omega_1$.
\item The model validates propositional extensionality, function
  extensionality, uniqueness of identity proofs, quotient types, and the
  $\Prop$-level Church's thesis, and refutes unique choice
  (Proposition~\ref{prop:validates}). Hence $\CIC_\omega+\EM\nvdash\uchoice$,
  and since $\CIC_2+\uchoice$ interprets $\IZF$ \cite{barras10} and
  $\IZF$ interprets $\ZF$ \cite{friedman73},
  $\CIC_\omega+\EM\nvdash\Con(\ZF)$ while $\CIC_2+\uchoice\vdash\Con(\ZF)$
  (Corollary~\ref{cor:boundary}).
\item $\CIC_2\vdash\Con(\HAH)$, and $\CIC_2\vdash\Con(\HAH_\beta)$ for every
  Brouwer ordinal term $\beta$ (Theorem~\ref{thm:lower}); by
  \cite{vdbergotten23}, $\CIC_1$ is conservative over $\HAH$.
\end{itemize}
The construction applies to Lean, which drops cumulativity, has general
inductive types, and adds definitional proof irrelevance and quotients, all
validated by the model (inductive types are handled as $W$-types are). So
$\ZF$ proves the consistency of Lean without \texttt{Classical.choice},
whereas Lean with choice proves $\Con(\ZFC+n\text{ inaccessibles})$ for every
$n$ \cite{carneiro19}: choice is exactly what the inaccessibles pay for.

\section{The type theories}\label{sec:theories}

$\CIC_n$ ($1\le n\le\omega$) is van den Berg and Otten's extensional theory
$\lambda C^+$ \cite[App.~A]{vdbergotten23} with its impredicative $\Set$
removed and a cumulative tower of predicative universes put in its place.
We give the rules in full, following their notation, since the exact type
formers matter for the model; the differences from $\lambda C^+$ are
collected in Remark~\ref{rem:diff}.

\paragraph{Syntax.} Sorts and terms are
\begin{align*}
s &::= \Prop \mid \Type_i \quad(i<n)\\
A,B,e &::= x \mid s \mid \Pi x{:}A.B \mid \lambda x{:}A.e \mid e\,e
  \mid \Sigma x{:}A.B \mid \langle e,e\rangle \mid \mathsf{ind}^\Sigma_C e
  \mid W x{:}A.B \mid \mathsf{tree}\,e\,e \mid \mathsf{ind}^W_C e\\
 &\quad\mid \mathbf n \mid \mathbf k_{\mathbf n} \mid \mathsf{ind}^{\mathbf n}_C e\cdots e
  \mid \N \mid 0 \mid \mathsf S\,e \mid \mathsf{ind}^\N_C e\,e
  \mid e =_A e \mid \mathsf{refl}\,e \mid \mathsf{ind}^=_C e\\
 &\quad\mid \|A\| \mid |e| \mid \mathsf{ind}^{\|\|}_C e
  \mid A/e \mid [e]_R \mid \mathsf{ind}^{/}_C e\,e \mid \mathsf{ax}_/
  \mid \mathsf{propext} \mid \mathsf{em}
\end{align*}
with $\mathbf n$ ranging over numerals and $\mathbf k_{\mathbf n}$ over
$k<n$. Contexts are $\Gamma::=\cdot\mid\Gamma,x{:}A$. Sorts are ordered
$\Prop<\Type_0<\Type_1<\cdots$ and $s_1\sqcup s_2$ is the maximum. The
judgements are $\Gamma\vdash a:A$ and $\Gamma\vdash a\equiv a':A$. We
abbreviate $A\to B$, $A\times B$ for non-dependent $\Pi$, $\Sigma$, and
\begin{gather*}
\bot:=\|\mathbf 0\|,\qquad
\neg P:=P\to\bot,\qquad
P\wedge Q:=\Sigma\_{:}P.Q,\qquad
P\leftrightarrow Q:=(P\to Q)\times(Q\to P),\\
P\vee Q:=\|\Sigma b{:}\mathbf 2.\ \mathsf{ind}^{\mathbf 2}_{\lambda\_.\Prop}\,P\,Q\,b\|,\qquad
\exists x{:}A.P:=\|\Sigma x{:}A.P\|.
\end{gather*}

\paragraph{Structural rules, sorts, cumulativity, products.}
In the rules $s,s_1,s_2$ range over sorts, and
$s_1\Rightarrow s_2:=\Prop$ if $s_2=\Prop$ and $s_1\sqcup s_2$ otherwise.
\[
\frac{\Gamma\vdash A:s}{\Gamma,x{:}A\vdash x:A}\ \text{\small$x\notin\Gamma$}\qquad
\frac{\Gamma\vdash A:s\quad\Gamma\vdash b:B}{\Gamma,x{:}A\vdash b:B}\ \text{\small$x\notin\Gamma$}\qquad
\frac{}{\vdash\Prop:\Type_0}\qquad
\frac{i+1<n}{\vdash\Type_i:\Type_{i+1}}
\]
\[
\frac{\Gamma\vdash A:s\quad s\le s'}{\Gamma\vdash A:s'}\qquad
\frac{\Gamma\vdash a:A\quad\Gamma\vdash A\equiv A':s}{\Gamma\vdash a:A'}\qquad
\frac{\Gamma\vdash p:a=_A a'}{\Gamma\vdash a\equiv a':A}\qquad
\frac{\Gamma\vdash P:\Prop\quad\Gamma\vdash p:P\quad\Gamma\vdash q:P}{\Gamma\vdash p\equiv q:P}
\]
\[
\frac{\Gamma\vdash A:s_1\quad\Gamma,x{:}A\vdash B:s_2}{\Gamma\vdash\Pi x{:}A.B:s_1\Rightarrow s_2}\qquad
\frac{\Gamma,x{:}A\vdash b:B}{\Gamma\vdash\lambda x{:}A.b:\Pi x{:}A.B}\qquad
\frac{\Gamma\vdash f:\Pi x{:}A.B\quad\Gamma\vdash a:A}{\Gamma\vdash f\,a:B[a/x]}\qquad
(\lambda x{:}A.b)\,a\equiv b[a/x]
\]
The rules stating that $\equiv$ is a congruence are omitted, as in
$\lambda C^+$; computation rules are written as unconditional equations
between well-typed terms. The first case of $\Rightarrow$ is impredicativity
of $\Prop$; for $n<\omega$ the top sort $\Type_{n-1}$ is not a type.

\paragraph{Formation.}
\[
\frac{\Gamma\vdash A:s_1\quad\Gamma,x{:}A\vdash B:s_2}{\Gamma\vdash\Sigma x{:}A.B:s_1\sqcup s_2}\qquad
\frac{\Gamma\vdash A:s_1\quad\Gamma,x{:}A\vdash B:s_2}{\Gamma\vdash W x{:}A.B:s_1\sqcup s_2}\qquad
\frac{}{\vdash\mathbf n:\Type_0}\qquad
\frac{}{\vdash\N:\Type_0}
\]
\[
\frac{\Gamma\vdash A:s\quad\Gamma\vdash a:A\quad\Gamma\vdash a':A}{\Gamma\vdash a=_A a':\Prop}\qquad
\frac{\Gamma\vdash A:s}{\Gamma\vdash\|A\|:\Prop}\qquad
\frac{\Gamma\vdash A:s\quad\Gamma,x{:}A,x'{:}A\vdash R:s'}{\Gamma\vdash A/R:s}
\]
So $\Prop$ is closed under $\Sigma$ and $W$ as well as $\Pi$; Coq's
singleton elimination and $\Acc$-recursion live here.

\paragraph{Introduction and elimination.}
Each rule is stated for a fixed well-formed type of the displayed shape;
$C$ is the motive, of any sort $s$ except where noted.
\[
\frac{\Gamma\vdash a:A\quad\Gamma\vdash b:B[a/x]}{\Gamma\vdash\langle a,b\rangle:\Sigma x{:}A.B}\qquad
\frac{\Gamma,p{:}\Sigma x{:}A.B\vdash C:s\quad\Gamma\vdash f:\Pi x{:}A.\Pi y{:}B.\,C[\langle x,y\rangle/p]}
     {\Gamma\vdash\mathsf{ind}^\Sigma_C f:\Pi p{:}\Sigma x{:}A.B.\,C}
\]
\[
\frac{\begin{array}{c}\Gamma,t{:}Wx{:}A.B\vdash C:s\\
      \Gamma\vdash f:\Pi a{:}A.\,\Pi d{:}(B[a/x]\to Wx{:}A.B).\,(\Pi b{:}B[a/x].\,C[d\,b/t])\to C[\mathsf{tree}\,a\,d/t]\end{array}}
     {\Gamma\vdash\mathsf{ind}^W_C f:\Pi t{:}Wx{:}A.B.\,C}
\]
\[
\frac{\Gamma\vdash a:A\quad\Gamma\vdash d:B[a/x]\to Wx{:}A.B}{\Gamma\vdash\mathsf{tree}\,a\,d:Wx{:}A.B}\qquad
\frac{k<n}{\vdash\mathbf k_{\mathbf n}:\mathbf n}\qquad
\frac{\Gamma,k{:}\mathbf n\vdash C:s\quad\Gamma\vdash c_j:C[\mathbf j_{\mathbf n}/k]\ (j<n)}
     {\Gamma\vdash\mathsf{ind}^{\mathbf n}_C c_0\cdots c_{n-1}:\Pi k{:}\mathbf n.\,C}
\]
\[
\frac{}{\vdash 0:\N}\qquad\frac{\Gamma\vdash m:\N}{\Gamma\vdash\mathsf S\,m:\N}\qquad
\frac{\Gamma,m{:}\N\vdash C:s\quad\Gamma\vdash c:C[0/m]\quad\Gamma\vdash f:\Pi m{:}\N.\,C\to C[\mathsf S\,m/m]}
     {\Gamma\vdash\mathsf{ind}^\N_C c\,f:\Pi m{:}\N.\,C}
\]
\[
\frac{\Gamma\vdash a:A}{\Gamma\vdash\mathsf{refl}\,a:a=_A a}\qquad
\frac{\Gamma,a{:}A,a'{:}A,e{:}a=_A a'\vdash C:s\quad\Gamma\vdash f:\Pi x{:}A.\,C[x,x,\mathsf{refl}\,x]}
     {\Gamma\vdash\mathsf{ind}^=_C f:\Pi a\,a'{:}A.\,\Pi e{:}a=_A a'.\,C}
\]
\[
\frac{\Gamma\vdash a:A}{\Gamma\vdash|a|:\|A\|}\qquad
\frac{\Gamma,t{:}\|A\|\vdash C:\Prop\quad\Gamma\vdash f:\Pi x{:}A.\,C[|x|/t]}
     {\Gamma\vdash\mathsf{ind}^{\|\|}_C f:\Pi t{:}\|A\|.\,C}\qquad
\frac{\Gamma\vdash a:A}{\Gamma\vdash[a]_R:A/R}
\]
\[
\frac{\begin{array}{c}\Gamma,q{:}A/R\vdash C:s\qquad\Gamma\vdash f:\Pi x{:}A.\,C[[x]_R/q]\\
      \Gamma\vdash h:\Pi x\,x'{:}A.\,R[x,x']\to f\,x=_{C[[x]_R/q]}f\,x'\end{array}}
     {\Gamma\vdash\mathsf{ind}^/_C f\,h:\Pi q{:}A/R.\,C}
\]
\[
\vdash\mathsf{ax}_/:\Pi a\,a'{:}A.\,R[a,a']\to[a]_R=_{A/R}[a']_R\qquad
\vdash\mathsf{propext}:\Pi P\,P'{:}\Prop.\,(P\leftrightarrow P')\to P=_\Prop P'
\]
\paragraph{Computation.}
\begin{gather*}
\mathsf{ind}^\Sigma_C f\langle a,b\rangle\equiv f\,a\,b\qquad
\mathsf{ind}^W_C f\,(\mathsf{tree}\,a\,d)\equiv f\,a\,d\,(\lambda b.\,\mathsf{ind}^W_C f\,(d\,b))\qquad
\mathsf{ind}^{\mathbf n}_C\vec c\ \mathbf j_{\mathbf n}\equiv c_j\\
\mathsf{ind}^\N_C c f\,0\equiv c\qquad
\mathsf{ind}^\N_C c f\,(\mathsf S\,m)\equiv f\,m\,(\mathsf{ind}^\N_C c f\,m)\qquad
\mathsf{ind}^=_C f\,a\,a\,(\mathsf{refl}\,a)\equiv f\,a\\
\mathsf{ind}^{\|\|}_C f\,|a|\equiv f\,a\qquad
\mathsf{ind}^/_C f\,h\,[a]_R\equiv f\,a
\end{gather*}
In the premise $h$ of $\mathsf{ind}^/$, the term $f\,x'$ has type
$C[[x']_R/q]$, which is convertible with $C[[x]_R/q]$ under the hypothesis
$R[x,x']$ by $\mathsf{ax}_/$ and reflection, so no transport is needed.
$\CIC_n+\EM$ adds $\vdash\mathsf{em}:\Pi P{:}\Prop.\,P\vee\neg P$.

\begin{remark}[differences from $\lambda C^+$]\label{rem:diff}
(i) $\lambda C^+$ has $\Prop,\Set:\Type$ with $\Set$ impredicative and
$\Type$ a top sort; we have $\Prop:\Type_0:\Type_1:\cdots$ with each
$\Type_i$ predicative, and $\CIC_1$ is $\lambda C^+$ minus $\Set$ (its
$\Type$ becomes our $\Type_0$).
(ii) Consequently $\Pi$, $\Sigma$, and $W$ land in the maximum of the sorts
involved rather than, as in $\lambda C^+$, in the sort of the codomain
(for $\Pi$) or of the domain (for $W$); those rules are sound for
impredicative sorts by \cite[Prop.~7]{vdbergotten23} but not for a
predicative universe.
(iii) $\lambda C^+$ allows $\mathsf{ind}^{\|\|}$ to eliminate into any
family $C$ of subsingletons, given $h:\Pi t.\,\Pi c\,c'{:}C.\,c=c'$. That
rule derives unique choice: in the context $R{:}\N\to\Prop,\
t{:}\|\Sigma m{:}\N.R\,m\|$, eliminate into the subsingleton
$C:=(\Sigma m{:}\N.R\,m)/\top$ with $f:=[{-}]$ to obtain, uniformly in $R$,
an element of $C$ whose realizer must encode some $m$ with $R\,m$. No such
tracker exists in the assembly model, where $\N\to\Prop$ is indiscrete,
and the clause of \cite[App.~B]{vdbergotten23} interpreting
$\mathsf{ind}^{\|\|}$ (``$f$ applied to any element of $A$'') does not
supply one. We therefore restrict the target of $\mathsf{ind}^{\|\|}$ to
$\Prop$, which is what the interpretation of $\HAH$ uses; see
Proposition~\ref{prop:validates} for the failure of unique choice.
(iv) Reflection makes the theory extensional, so funext and UIP are
derivable. Intensional variants (Coq's $J$ with or without UIP and funext,
Lean's definitional proof irrelevance and $\mathsf{Quot}$) are subsystems
up to notation, and the model validates the extensional rules, so nothing
below depends on the choice. General inductive types are deliberately
absent, as in $\lambda C^+$; $\N$, $\mathbf 2$, $\Sigma$, and $W$ suffice
for this note.
\end{remark}

\section{Lower bounds}\label{sec:lower}

The reason one universe is exactly $\HAH$ and two are more is that a universe
$\Type_0:\Type_1$ lets us define types by recursion on $\N$.

\begin{theorem}\label{thm:lower}
$\CIC_2\vdash\Con(\HAH)$. More generally, for each closed term $\beta$ of the
type $O$ of Brouwer ordinals ($O:=\mathsf{zer}\mid\mathsf{suc}\,O\mid
\mathsf{lim}\,(\N\to O)$), $\CIC_2$ proves the consistency of
$\HAH_\beta$, higher-order Heyting arithmetic with a sort for every ordinal
level below $\beta$, where a limit level $\mathsf{lim}\,f$ is the sum
$\Sigma(n:\N)\,\mathsf{lvl}(f\,n)$.
\end{theorem}
\begin{proof}
Let $O:=W x{:}\N.\,B\,x:\Type_0$ with $B\,0=\mathbf 0$, $B\,1=\mathbf 1$,
$B\,(m{+}2)=\N$ defined by $\mathsf{ind}^\N$ into $\Type_0$, which is a
type only because $\Type_0:\Type_1$. Define $T:O\to\Type_0$ by
$\mathsf{ind}^W$, again into $\Type_0$:
$T(\mathsf{zer})=\N$, $T(\mathsf{suc}\,o)=T(o)\to\Prop$,
$T(\mathsf{lim}\,f)=\Sigma m{:}\N.\,T(f\,m)$. Each $T(o):\Type_0$ since
$\Prop$ is impredicative and $\Prop:\Type_0$. Code formulas of $\HAH_\beta$
as a second $W$-type over $\N$ (labels for atoms, connectives, and
quantifiers with their level and variable), environments as
$\Pi o{:}O.\,\N\to T(o)$, and satisfaction
$\mathsf{Sat}:\mathsf{form}\to\mathsf{Env}\to\Prop$ by $\mathsf{ind}^W$ with
motive $\mathsf{Env}\to\Prop:\Type_0$, so that
$\mathsf{Sat}(\forall^o x.A)\,\rho=\Pi v{:}T(o).\,\mathsf{Sat}(A)(\rho[x\mapsto v])$.
Every axiom of $\HAH_\beta$ is satisfied (comprehension at level $o+1$ is
$\lambda$-abstraction into $\Prop$), so by induction on derivations no
proof of $\bot$ exists, which is a $\Pi^0_1$ statement about codes of
derivations.
\end{proof}

In $\CIC_1$ the motives $\N\to\Type_0$ and $O\to\Type_0$ are not types, and indeed
$\lambda C^+$ is conservative over $\HAH$ \cite{vdbergotten23}. The rest of
this note shows that the climb from $\CIC_1$ to $\CIC_\omega$ nevertheless
stays far below $\ZF$.

\section{The model}\label{sec:model}

\subsection{Assemblies and a classical \texorpdfstring{$\Prop$}{Prop}}

We use Kleene application $e\ap k$ on $\N$, with $e\ap k{\downarrow}$ for
definedness. An \emph{assembly} $A$ is a set $|A|$ with a relation
${\rz_A}\subseteq\N\times|A|$ such that every $x\in|A|$ has some $k\rz_A x$;
a morphism $A\to B$ is a function $F$ tracked by some $e$, meaning
$e\ap k{\downarrow}\rz_B F(x)$ whenever $k\rz_A x$. $A$ is \emph{modest} if
distinct elements have disjoint realizer sets, so that a realizer determines
the element and a morphism out of a modest set is determined by its tracker.
$\Sigma$, $\Pi$, $W$ of
assembly-valued families are as in \cite[\S5.3]{vdbergotten23}: the carrier
is the set of set-theoretic elements that have a realizer, with $\Pi$
realized by trackers, $\Sigma$ by pairs, and $W$ (elements represented as
sets of finite paths) by programs producing subtree realizers. A quotient
$A/R$ by a relation $R\subseteq|A|^2$ has carrier $|A|/{\sim_R}$ (the
equivalence relation generated) and $k\rz[x]$ iff $k\rz_A y$ for some
$y\sim_R x$. $\mathbf n$ and $\N$ are the evident modest assemblies. Unlike \cite{vdbergotten23} we do not confine assemblies to
$\PP^n(\N)$: $\Asm$ is the class of all assemblies in $\ZF$.

We take $\Prop$ classical: $|\Prop|=\{\bot,\top\}$ with
$\rz_\Prop=\N\times|\Prop|$. As a type, $\bot$ decodes to $\mathbf 0$ and
$\top$ to $\mathbf 1$. The connectives are Boolean: $\Pi(x{:}A)\,p(x)$ is
$\top$ iff every $p(x)$ is, for \emph{any} assembly $A$ (impredicativity);
$\Sigma(x{:}p)\,q$ is $p\wedge q$; $W(x{:}p)\,q$ is $p\wedge\neg q$ (a
well-founded tree exists iff the root label exists and it has no children);
$[x=_A y]$ is equality in $|A|$; $\|A\|$ is $[\,|A|\ne\emptyset\,]$; and a
quotient of a proposition is itself. This is Werner's proof-irrelevant
$\Prop$ \cite{werner97,miquelwerner02} sitting on assemblies instead of
sets.

\subsection{The universe operator}

A universe will be a set of \emph{codes}. Codes are not natural numbers:
a type such as $\Sigma x{:}\N.\,P\,x$ depends on an arbitrary subset
$P\subseteq\N$, so a universe containing $\N$ and $\Prop$ has at least
continuum many elements. Codes are instead well-founded trees whose nodes
carry the families themselves, in the style of Aczel's and Werner's set
universes. What keeps the construction small is not the number of codes but
the number of \emph{stages}: a code is realized by a natural number that
determines its computational skeleton, and the stage at which a code enters
the universe depends only on that skeleton.

\begin{definition}[codes, pre-universes, families]\label{def:codes}
The class of codes is defined inductively: $\bot,\top,\code\Prop,
\code{\mathbf n}\ (n\in\N),\code\N,\code{\U_i}$ are codes; and if $c$ is a code,
$X$ a set, $F:X\to\text{codes}$, and $R\subseteq X^2$, then
$\code\Pi(c,F)$, $\code\Sigma(c,F)$, $\code W(c,F)$, and $\code{/}(c,R)$
are codes, \emph{except} that the propositional cases are identified with
truth values: $\code\Pi(c,F):=[\forall x.\,F(x)=\top]$ if every $F(x)$ is
a truth value; $\code\Sigma(c,F):=[c=\top\wedge F(\star)=\top]$ and
$\code W(c,F):=[c=\top\wedge F(\star)=\bot]$ if $c$ and every $F(x)$ are
truth values; and $\code/(c,R):=c$ if $c$ is a truth value. (Here $X$ will
always be the carrier of the decoding of $c$.)
A \emph{pre-universe} is a set $u$ of pairs (code, assembly) that is a
partial function, written $u(c)$. Realizers of codes are defined by
recursion on codes: every $k$ realizes a truth value or a constant code;
$\langle k,e\rangle$ realizes $\code Q(c,F)$ iff $k$ realizes $c$ and $e$
tracks $F$ as a map $u(c)\to\langle u\rangle$; $k$ realizes $\code/(c,R)$
iff $k$ realizes $c$. The \emph{universe assembly} $\langle u\rangle$ has
carrier $\dom u$ and this realizability relation. A \emph{$u$-family over
$c\in\dom u$} is a morphism $F:u(c)\to\langle u\rangle$ in $\Asm$, and
$u\circ F$ is the assembly-valued family $x\mapsto u(F(x))$.
\end{definition}

The identification of propositional codes with truth values is the
code-level form of Miquel and Werner's trick \cite{miquelwerner02}: it makes
$\Prop\subseteq\Type_0$ hold on the nose and makes the interpretation of a
type independent of the sort at which it is formed.

\begin{definition}[the closure operator]
For pre-universes $v$ (the \emph{seed}) and $u$ let
\begin{align*}
\Phi_v(u) \;=\;& u\;\cup\;v\;\cup\;
  \{(\bot,\mathbf 0),\ (\top,\mathbf 1),\ (\code\Prop,\Prop),\ (\code\N,\N)\}
  \cup\{(\code{\mathbf n},\mathbf n):n\in\N\}\\
&\cup\;\{(\code{Q}(c,F),\ Q(u(c),u\circ F)) : Q\in\{\Pi,\Sigma,W\},\
      c\in\dom u,\ F\text{ a $u$-family over }c\}\\
&\cup\;\{(\code/(c,R),\ u(c)/R) : c\in\dom u,\ R\subseteq|u(c)|^2\}.
\end{align*}
Let $\El_0:=\mathrm{lfp}\,\Phi_\emptyset$ and
$\El_{n+1}:=\mathrm{lfp}\,\Phi_{v_n}$ with seed
$v_n:=\El_n\cup\{(\code{\U_n},\langle\El_n\rangle)\}$, where
$\mathrm{lfp}\,\Phi_v:=\bigcup_\alpha u^\alpha$ for the stages
$u^\alpha:=\Phi_v(\bigcup_{\beta<\alpha}u^\beta)$.
\end{definition}

For the propositional cases the clauses agree with the truth-value
definitions of \S4.1, so the operator is well defined on codes that have
been identified: e.g.\ $\Pi(u(c),u\circ F)$ with all $F(x)$ truth values is
$\mathbf 1$ or $\mathbf 0$ according as $[\forall x.\,F(x)=\top]$, and a
$W$ over $\mathbf 1$ with children indexed by $\mathbf 1$ is empty.

\begin{lemma}\label{lem:lfp}
For every $n<\omega$:
\begin{enumerate}
\item $\Phi_{v_n}$ is monotone with respect to $\subseteq$, and each stage
  $u^\alpha$ is a pre-universe extending the previous stages.
\item The stages stabilise at a countable ordinal $\sigma_n$, so $\El_n$ is
  a set, and it is the least $u\supseteq v$ with $\Phi_v(u)=u$ for its seed
  $v$ ($v=\emptyset$ or $v=v_{n-1}$).
\item $\dom\El_n\subseteq\dom\El_{n+1}$ and $\El_{n+1}$ restricted to
  $\dom\El_n$ is $\El_n$; moreover $\El_{n+1}(\code{\U_n})=\langle\El_n\rangle$
  and $\Prop$ is a sub-assembly of $\langle\El_0\rangle$.
\item Every value $\El_n(c)$ is an assembly of rank below $\omega_1$, so
  $\langle\El_n:n<\omega\rangle\in V_{\omega_1+1}$; the carriers, however,
  have sizes cofinal in $\beth_{\sigma_n}$.
\end{enumerate}
\end{lemma}
\begin{proof}
(1) Every clause of $\Phi_v$ is positive in $u$: enlarging $u$ enlarges
$\dom u$, hence the set of $u$-families, and leaves the values $u(c)$ and
the realizers of existing codes unchanged. A code is its own syntax tree,
so a code can only be assigned one value, and a code already present was
added from a smaller stage whose restriction agrees with the current one.
(2) Define the \emph{stage} of a code $c\in\dom\El_n$ as the least $\alpha$
with $c\in u^\alpha$. We claim a realizer of $c$ determines its stage. For
constants and truth values the stage is $0$ or the seed's. For
$\code Q(c,F)$ with realizer $\langle k,e\rangle$: by induction $k$
determines the stage of $c$; and for each $x$ and $k'\rz x$, $e\ap k'$
realizes $F(x)$ and so determines its stage, hence the countable set
$\{e\ap k':k'\in\N\}$ determines the set of stages of the $F(x)$; the stage
of the code is one above the supremum of all these. Quotients are one above
the stage of $c$. Now if $u^{\alpha+1}\ne u^\alpha$, the least realizer of a
code entering at $\alpha+1$ is a natural number not assigned to any other
stage; this injects the strict stages into $\N$, so $\sigma_n<\omega_1$,
with no use of choice. Each stage is a set by Replacement.
(3) $v_n\subseteq\El_{n+1}$ by the seed clause, and the closure clauses
applied to codes of $\El_n$ give the same values as inside $\El_n$. The
code $\code{\U_n}$ is produced by no earlier clause. The truth values are
in $\dom\El_0$ with the indiscrete realizers, which is exactly $\Prop$.
(4) By induction on the stage: the constants and $\langle\El_n\rangle$ have
rank $<\omega+\omega$, and each type former raises the supremum of the
ranks of its arguments by a finite amount, so $\rk\El_n(c)<\omega+\omega
\cdot(\sigma_n+1)<\omega_1$. Sizes grow because $\Sigma x{:}A.\,P\,x$ for
$P\subseteq|A|$ gives $2^{|A|}$ distinct types.
\end{proof}

The heart of the matter is (2). In the naive model a family over
$A\in V_\kappa$ is any function $|A|\to V_\kappa$, and the stage at which
$\Pi(A,F)$ appears is one above the supremum of the stages of the $F(x)$,
which is cofinal in $\kappa$ unless $\kappa$ is regular; closure under
powerset then makes it a strong limit. Here the stages of the $F(x)$ are
those of the countably many values $e\ap k'$ of the tracker, so the
supremum of a countable set of countable ordinals suffices, and the
closure stops at a countable stage. The universes are still enormous, but
size is harmless: each $\Pi$, $\Sigma$, $W$, and quotient only has to
\emph{exist}, which $\ZF$ grants, and nothing has to be closed under all
functions.

The preceding paragraph and Lemma~\ref{lem:lfp}(2) are wrong, and not
repairably so. The claim that a realizer determines the stage of a code
fails: $\code\Pi(\code\Sigma(\code\N,P),\ m\mapsto d_m)$ with $d_m$ a code
of stage $m$ is realized, for every $P\subseteq\N$, by the same pair
$\langle\langle k_\N,0\rangle,e_d\rangle$, but has stage $3$ for
$P=\{0\}$ and stage $\omega+1$ for $P=\N$. Worse, the closure never
stabilises.

\begin{proposition}\label{prop:noset}
Assume $\ZFC$. Let $U$ be a set of codes containing $\code\N$ and
$\code\Prop$ and closed under $\Pi$ and $\Sigma$ over tracked families as
in Definition~\ref{def:codes}. Then $\kappa:=\sup\{|\El(c)|:c\in U\}$ is a
strongly inaccessible cardinal. Hence, in the absence of inaccessibles,
$\mathrm{lfp}\,\Phi_\emptyset$ is a proper class and there is no assembly
model of $\CIC_2$ in which $\Type_0$ is an assembly.
\end{proposition}
\begin{proof}
For $c\in U$ the code $\code\Pi(c,x\mapsto\code\Prop)$ decodes to the full
function space $|\El c|\to\{\bot,\top\}$, because $\Prop$ is indiscrete,
so $2^{|\El c|}$ is a size occurring in $U$; hence every size is $<\kappa$
and $\kappa$ is a strong limit, uncountable since $\PP(\N)$ occurs. Every
code has a realizer, so with $\mu_n:=\sup\{|\El c|:n\rz c\}$ we have
$\kappa=\sup_n\mu_n$. If $\mu_n<\kappa$ for all $n$, choose
$c_n$ with $n\rz c_n$ and $\sup_n|\El c_n|=\kappa$; the family
$n\mapsto c_n$ is a morphism $\N\to\langle U\rangle$ tracked by the
identity, so $\code\Sigma(\code\N,n\mapsto c_n)\in U$ has size $\kappa$,
contradicting the first sentence. So $\mu_{n_0}=\kappa$ for some $n_0$.
Now let $A\in U$ and let $F:|\El A|\to\{c:n_0\rz c\}$ be \emph{any}
function; it is tracked by the constant $n_0$, so
$\code\Sigma(A,F)\in U$ and $\sum_x|\El F(x)|<\kappa$. Choosing $F$ with
image cofinal in $\kappa$ whenever $|\El A|\ge\operatorname{cf}\kappa$
would contradict this, so $|\El A|<\operatorname{cf}\kappa$ for every
$A\in U$, i.e.\ $\operatorname{cf}\kappa=\kappa$.
\end{proof}

The two ingredients are exactly the ones the type theory forces: $\Prop$
must be indiscrete for the impredicative $\forall$ to be tracked, which
makes $A\to\Prop$ the set-theoretic powerset; and the context
$X{:}\Type_0,\ f{:}X\to\Type_0\vdash\Sigma x{:}X.\,f\,x:\Type_0$ forces
closure under $\Sigma$ over \emph{every} morphism into the universe. In
$\lambda C^+$ neither bites, because its only set-sized universes are
$\mathrm{Subsing}$ and $\mathrm{PER}$, which are closed under products
over arbitrary assemblies for the reason recorded in
Remark~\ref{rem:diff}(ii), and $\Type$ is a class. The trouble starts
precisely with $\Prop:\Type_0:\Type_1$, the same configuration that gives
the lower bound of Theorem~\ref{thm:lower}. So the inaccessibles of
\cite{werner97,barras10,timanysozeau17} are not an artifact of $V_\kappa$:
they are forced in every tracked model with a universe object. Whether
they are forced proof-theoretically, i.e.\ whether $\ZF\vdash\Con(\CIC_2)$,
is open; note the symmetry that unique choice is what the model side
lacks to keep universes small (Proposition~\ref{prop:validates}) and what
the syntactic side lacks to prove Replacement \cite{barras10}.

\subsection{The interpretation}

For each sort $s$ let $\U_s$ be its universe object and $\El_s$ its
decoding: $\U_\Prop:=\Prop$ with $\El_\Prop(\bot)=\mathbf 0$,
$\El_\Prop(\top)=\mathbf 1$; $\U_{\Type_i}:=\langle\El_i\rangle$ with
$\El_{\Type_i}:=\El_i$; and for the top sort of $\CIC_n$ ($n<\omega$),
$\U_{\Type_{n-1}}:=\Asm$ with $\El$ the identity. By
Lemma~\ref{lem:lfp}(3), $\U_s\subseteq\U_{s'}$ and the decodings agree
whenever $s\le s'$. Following \cite[App.~B]{vdbergotten23}, by simultaneous
induction on derivations we define
\begin{itemize}
\item for each context $\Gamma$ an assembly $\sem\Gamma$;
\item for each $\Gamma\vdash A:s$ a morphism $\sem{\Gamma\vdash A:s}:\sem\Gamma\to\U_s$
  (for the top sort: a family of assemblies indexed by $|\sem\Gamma|$);
  we write $\sem{\Gamma\vdash A}(G)$ for $\El_s(\sem{\Gamma\vdash A:s}(G))$;
\item for each $\Gamma\vdash a:A$ a section $\sem{\Gamma\vdash a:A}$, i.e.\ a
  tracked function $G\mapsto\sem{\Gamma\vdash a:A}(G)\in|\sem{\Gamma\vdash A}(G)|$.
\end{itemize}
The identifications of Definition~\ref{def:codes} make the clauses depend
only on the conclusion of the rule, so the interpretation is independent of
the derivation; in particular cumulativity, conversion, and the sort at
which a type is formed do not change its code.
\begin{gather*}
\sem{\cdot}:=\mathbf 1,\qquad
\sem{\Gamma,x{:}A}:=\Sigma(G\in\sem\Gamma)\,\sem{\Gamma\vdash A}(G),\\
\sem{\Gamma,x{:}A\vdash x:A}(\langle G,X\rangle):=X,\qquad
\sem{\Gamma,x{:}A\vdash b:B}(\langle G,X\rangle):=\sem{\Gamma\vdash b:B}(G),\\
\sem{\Gamma\vdash a:A'}:=\sem{\Gamma\vdash a:A}\ \text{(conversion)},\qquad
\sem{\Gamma\vdash A:s'}:=\sem{\Gamma\vdash A:s}\ \text{(cumulativity)},\\
\sem{\vdash\Prop:\Type_0}:=\code\Prop,\qquad
\sem{\vdash\Type_i:\Type_{i+1}}:=\code{\U_i},\qquad
\sem{\vdash\mathbf n:\Type_0}:=\code{\mathbf n},\qquad
\sem{\vdash\N:\Type_0}:=\code\N,\\
\sem{\vdash\mathbf k_{\mathbf n}:\mathbf n}:=k,\qquad
\sem{\Gamma\vdash\mathsf{ind}^{\mathbf n}_C\vec c:\Pi k{:}\mathbf n.C}(G)(k):=\sem{\Gamma\vdash c_k:C[\mathbf k_{\mathbf n}/k]}(G),\\
\sem{\vdash 0:\N}:=0,\qquad
\sem{\Gamma\vdash\mathsf S\,m:\N}(G):=\sem{\Gamma\vdash m:\N}(G)+1,\\
\sem{\Gamma\vdash\mathsf{ind}^\N_C c f:\Pi m{:}\N.C}(G)(0):=\sem{\Gamma\vdash c:C[0/m]}(G),\\
\sem{\Gamma\vdash\mathsf{ind}^\N_C c f:\Pi m{:}\N.C}(G)(m{+}1):=\sem{\Gamma\vdash f}(G)(m)\bigl(\sem{\Gamma\vdash\mathsf{ind}^\N_C cf}(G)(m)\bigr).
\end{gather*}
For the type formers with $Q\in\{\Pi,\Sigma,W\}$, and for quotients,
\begin{align*}
\sem{\Gamma\vdash Qx{:}A.B:s}(G)&:=\code Q\bigl(\sem{\Gamma\vdash A:s_1}(G),\ X\mapsto\sem{\Gamma,x{:}A\vdash B:s_2}(\langle G,X\rangle)\bigr),\\
\sem{\Gamma\vdash A/R:s}(G)&:=\code/\bigl(\sem{\Gamma\vdash A:s}(G),\ \{(X,X'):\sem{\Gamma,x,x'\vdash R}(\langle G,X,X'\rangle)\ne\emptyset\}\bigr),
\end{align*}
which are in $\dom\El_s$ by the closure clauses, since the second
component is a morphism into $\U_s$ tracked by the tracker of the family
$\sem{\Gamma,x{:}A\vdash B:s_2}$ composed with pairing; when $s=\Prop$ the
code is a truth value and this is the Boolean clause of \S4.1. The term
clauses are those of \cite[App.~B]{vdbergotten23}: $\lambda$ and
application are currying and evaluation, $\langle a,b\rangle$ is the pair,
$\mathsf{ind}^\Sigma_C f\,\langle X,Y\rangle:=\sem f(G)(X)(Y)$,
$\mathsf{tree}\,a\,d$ is the tree with root $\sem a(G)$ and subtrees
$\sem d(G)(Y)$, $\mathsf{ind}^W_C f\,T:=\sem f(G)(\mathrm{root}\,T)
(\text{subtrees})(Y\mapsto\mathsf{ind}^W_C f\,(\text{subtree}_Y))$,
$[a]_R$ is the class of $\sem a(G)$, and $\mathsf{ind}^/_C f\,h\,[X]:=
\sem f(G)(X)$. For the propositional constructs,
\begin{align*}
\sem{\Gamma\vdash a=_A a':\Prop}(G)&:=[\sem a(G)=\sem{a'}(G)], &
\sem{\Gamma\vdash\|A\|:\Prop}(G)&:=[\,|\sem{\Gamma\vdash A}(G)|\ne\emptyset\,],\\
\sem{\mathsf{refl}\,a}(G)&:=\star, &
\sem{\mathsf{ind}^=_C f}(G)(X)(X)(\star)&:=\sem f(G)(X),\\
\sem{|a|}(G)&:=\star, &
\sem{\mathsf{ind}^{\|\|}_C f}(G)(\star)&:=\star,
\end{align*}
and $\sem{\mathsf{ax}_/}=\sem{\mathsf{propext}}=\sem{\mathsf{em}}:=\star$,
which is well typed because the corresponding propositions are $\top$ by
the definitions of quotients, of $\Prop$, and of $\vee$.

\begin{theorem}\label{thm:main}
$\ZF$ proves that the above is a model of $\CIC_\omega+\EM$; hence
$\ZF\vdash\Con(\CIC_\omega+\EM)$.
\end{theorem}
\begin{proof}[Proof sketch]
One checks that each clause is well defined, tracked, and validates the
computation rules and reflection. Well-definedness of the type clauses is
Lemma~\ref{lem:lfp}; tracking of the type clauses is the pairing of
trackers noted above; the term clauses are tracked exactly as in
\cite[App.~B]{vdbergotten23}, with $\mathsf{ind}^\N$ and $\mathsf{ind}^W$
tracked by primitive recursion and a fixed-point combinator that terminates
on realizers of a well-founded tree by induction on the tree. When the
motive is a family of codes, e.g.\ $T$ in Theorem~\ref{thm:lower}, the
same program tracks a family into $\langle\El_i\rangle$, since realizers of
codes are computed structurally. $\mathsf{ind}^/$ is tracked by the tracker
of $f$ because realizers of a class are realizers of its representatives,
and $h$ together with set-theoretic equality makes $\sem f(G)(X)$
independent of $X$; $\mathsf{ind}^{\|\|}$ into $\Prop$ needs no tracker.
Reflection holds because $[\sem a(G)=\sem{a'}(G)]=\top$ forces
$\sem a=\sem{a'}$, proof irrelevance because every proposition decodes to
$\mathbf 0$ or $\mathbf 1$, and $\EM$ because $|\Prop|$ has two elements.
Elimination out of a propositional $\Sigma$ or $W$ into $\Type_i$ is sound
because their decodings are $\mathbf 0$ or $\mathbf 1$ and the minor
premise is realized whenever the type is inhabited. Finally
$\bot=\|\mathbf 0\|$ decodes to $\mathbf 0$, so no closed term of type
$\bot$ is interpretable. The construction of $\langle\El_n\rangle_{n<\omega}$
is a single set by Lemma~\ref{lem:lfp}, so the induction is uniform in $n$
and $\Con(\CIC_\omega+\EM)$ is one $\Pi^0_1$ statement.
\end{proof}

\begin{remark}[where the set theory is used]
Zermelo set theory is not enough: the stages and ranks in
Lemma~\ref{lem:lfp} climb along $\sigma_n$, which is countable but
recursively large (already $\El_0$ contains $\code{T(\beta)}$ for every
recursive Brouwer ordinal $\beta$, as $\beta\mapsto\code{T(\beta)}$ is
computable), and the carriers climb through $\beth_{\sigma_n}$. Beyond that
only $V_{\omega_1}$ is needed, and no choice: the injection in
Lemma~\ref{lem:lfp}(2) is canonical and Kleene application is definable.
\end{remark}

\begin{remark}[other choices of $\Prop$]\label{rem:prop}
Nothing in Lemma~\ref{lem:lfp} depends on what $\Prop$ is: propositions
are leaves of codes, of stage $0$, and only the truth-value clauses of
Definition~\ref{def:codes} change. Four choices realise every combination
of $\EM$ and $\mathsf{propext}$.
\begin{itemize}
\item $|\Prop|=\{\bot,\top\}$, as above: $\EM$ and $\mathsf{propext}$.
\item $|\Prop|=\PP(\N)$, indiscrete, with $P$ true iff $P\ne\emptyset$,
  connectives any Boolean-correct ones, and $=_\Prop$ set equality: $\EM$
  holds but $\mathsf{propext}$ fails ($\N$ and $\{0\}$ are both true).
\item $|\Prop|=\PP(\N)$ with realizability connectives
  ($P\to Q=\{e:\forall k\in P.\,e\ap k\in Q\}$, $\forall$ an intersection,
  $\vee$ by tagging) and $=_\Prop$ set equality: neither $\EM$ (no uniform
  decider) nor $\mathsf{propext}$.
\item Realizability connectives with $=_\Prop$ interpreted as realizable
  interderivability $P\leftrightarrow Q$: $\mathsf{propext}$ holds and $\EM$
  fails. This is the intuitionistic variant of Section~\ref{sec:discussion},
  and it requires $\Prop$ to carry an equality other than set equality (a
  setoid or PER-style object) so that reflection can be validated; one
  cannot instead quotient $\PP(\N)$ by interderivability, since any two
  nonempty realizer sets are interderivable via a constant program and the
  quotient is the first option. Van den Berg and Otten sidestep this by
  working in $\HAH$, where $\PP(\{\star\})$ is not provably two-valued;
  read classically, their $\mathrm{Subsing}$ is our first option.
\end{itemize}
In every case unique choice fails, by the argument of
Proposition~\ref{prop:validates}.
\end{remark}

\section{What the model decides}\label{sec:decides}

\begin{proposition}\label{prop:validates}
The model of Theorem~\ref{thm:main} validates propositional extensionality,
function extensionality, uniqueness of identity proofs, proof irrelevance,
quotient types with $\mathsf{Quot.sound}$ and $\mathsf{Quot.lift}$, and the
$\Prop$-level Church's thesis
$\CT^\Prop:=\forall f{:}\N\to\N,\;\exists e,\;\forall n,\;f\,n=e\ap n$.
It refutes the unique choice axiom
\[
\uchoice:\quad \forall (A\,B{:}\Type)\,(R{:}A\to B\to\Prop),\
(\forall x,\exists! y,R\,x\,y)\to\exists f{:}A\to B,\ \forall x,\ R\,x\,(f\,x).
\]
\end{proposition}
\begin{proof}
Propositions are truth values (propext, proof irrelevance), morphisms are
functions (funext), equality is equality (UIP), and a quotient is the
assembly of equivalence classes with the union of realizers, with
$\mathsf{Quot.lift}$ tracked by the tracker of the lifted function. Every
element of the assembly $\N\to\N$ is computable, and $\exists$ in $\Prop$
demands no realizer, so $\CT^\Prop$ is $\code\top$. For $\uchoice$ take
$A=B=\N$ and $R\,x\,y:=(y=1\leftrightarrow x\in K)$ with $K$ the halting set
(a $\Sigma^0_1$ formula): $\forall x,\exists!y,R\,x\,y$ is $\code\top$, but
the only candidate $f$ is the characteristic function of $K$, which has no
tracker, so $\exists f$ is $\code\bot$.
\end{proof}

\begin{corollary}\label{cor:boundary}
$\CIC_\omega+\EM\nvdash\uchoice$ and $\CIC_\omega+\EM\nvdash\Con(\ZF)$,
whereas $\CIC_2+\uchoice\vdash\Con(\ZF)$. Consequently $\ZF$ proves
$\Con(\CIC_\omega+\EM)$ but not $\Con(\CIC_2+\uchoice)$.
\end{corollary}
\begin{proof}
The first claim is Proposition~\ref{prop:validates}. For the second, if
$\CIC_\omega+\EM\vdash\Con(\ZF)$ then, arguing in $\ZF$, $\neg\Con(\ZF)$
would give a $\ZF$-proof of $\bot$, and by $\Sigma^0_1$-completeness
$\CIC_\omega+\EM$ would prove $\neg\Con(\ZF)$ as well and be inconsistent,
contradicting Theorem~\ref{thm:main}; so $\ZF\vdash\Con(\ZF)$, which is
impossible. For the third, Barras's sets-as-trees interpretation
\cite{barras10} of $\IZF$ in Coq with $\uchoice$ uses $\Type_0:\Type_1$
only, and Friedman's interpretation of $\ZF$ in $\IZF$ \cite{friedman73}
is arithmetically formalisable, so $\CIC_2+\uchoice\vdash\Con(\IZF)\to
\Con(\ZF)$. The last claim follows from the second and third.
\end{proof}

So the answer to ``does $\CIC$ prove $\Con(\ZF)$?'' is no, even with excluded
middle, and the reason is precisely the principle that Werner and Barras had
to add: functional or unique choice, which the assembly model refutes. In Lean
this principle is a consequence of \texttt{Classical.choice}, so the same
corollary reads: Lean minus choice is consistent relative to $\ZF$, and choice
is what the $\omega$ inaccessibles of \cite{carneiro19} are paying for.

\section{Discussion}\label{sec:discussion}

\paragraph{The intuitionistic variant.}
Taking instead $|\Prop|=\PP(\N)$ with $\forall(x{:}A)\,p(x)=\bigcap_x p(x)$
(uniform realizers) gives the model of \cite{vdbergotten23} extended by our
universes. Propext then holds with equality of propositions being realizable
interderivability and $\EM$ fails; the truth values in
Definition~\ref{def:codes} become realizer sets and the skeleton argument
of Lemma~\ref{lem:lfp}(2) is unchanged. One must
\emph{not} quotient $\PP(\N)$ by interderivability: any two nonempty
realizer sets are interderivable via a constant program, so the quotient is
the two-valued $\Prop$ above. The intuitionistic variant is what an
extension of the conservativity theorem of \cite{vdbergotten23} to $n$
universes would need; the natural conjecture is conservativity over
$\HAH_\sigma$ for $\sigma$ tied to the closure ordinals $\sigma_n$, which we
have not computed.

\paragraph{Impredicative $\Set$.}
$\lambda C^+$ has a second impredicative sort $\Set$ carrying data. The
classical collapse is unavailable there, since impredicative $\Set$ with an
informative excluded middle is inconsistent \cite{chicli02}, and a countable
Tarski universe is not impredicative. Van den Berg and Otten interpret
$\Set$ as $\nabla\mathrm{PER}$, the modest sets, which is closed under
products over arbitrary assemblies because morphisms into a modest set are
determined by their trackers; predicative universes above it should go
through as here, with $\nabla\mathrm{PER}$ as one more constant clause.

\paragraph{Induction-recursion and Agda.}
The construction is itself an inductive-recursive definition (codes with a
decoding) carried out as a least fixed point in $\ZF$, and we expect the
same countability argument to model a universe closed under
inductive-recursive definitions as in Agda. Dybjer and Setzer's naive model
needs a Mahlo cardinal \cite{dybjersetzer99} for the same reason the naive
$\CIC$ model needs inaccessibles: it closes under \emph{all} set-theoretic
operators, whereas the type theory can name only countably many.

\paragraph{What remains to be checked.}
The new ingredient is Lemma~\ref{lem:lfp}; the rest of
Theorem~\ref{thm:main} is \cite{werner97,miquelwerner02,leewerner11,
vdbergotten23} with $V_\kappa$ replaced by $\langle\El_i\rangle$. Not
written out here are (a) the verification that the identifications in
Definition~\ref{def:codes} commute with every clause of $\Phi$ (we checked
$\Pi$, $\Sigma$, $W$, and quotients) and (b) the tracking of the recursors
when the motive is a family of codes. A complete presentation would be a categories-with-families
model in which each $\langle\El_i\rangle$ is a universe object, a natural
target for formalisation.

{\small
\bibliographystyle{plain}
\bibliography{references}}

\end{document}
